% answers_short.tex —— 2006 年·普通高等学校招生考试·福建卷理科 · 简短答案速查
% 由 tools/make_short_answers.py 自 answers.tex 的 \gksolhead/\gkanswer 提取，
% 并与 2006/2006_数学真题_福建理_答案要点.md 逐题比对后生成，勿手改（改源 MD 后重跑）。

\begin{gknotice}
\textbf{使用说明}：本节只给最终答案便于快速核对，完整推导见解析版。答案为编写者独立推导，\textbf{非官方标准答案}。
\end{gknotice}

\gkpart{选择题}
\begin{center}\begin{tabular}{ccccccccccccc}
\toprule
题号 & 1 & 2 & 3 & 4 & 5 & 6 & 7 & 8 & 9 & 10 & 11 & 12 \\
\midrule
答案 & D & B & A & C & D & A & C & A & B & C & B & A \\
\bottomrule
\end{tabular}\end{center}

\gkpart{填空题}
\begin{itemize}[leftmargin=2.6em]
  \item[\textbf{13.}] $10$
  \item[\textbf{14.}] $\dfrac14$
  \item[\textbf{15.}] $\dfrac49$
  \item[\textbf{16.}] $\left(\dfrac53,\dfrac23\right)$
\end{itemize}

\gkpart{解答题}
\begin{itemize}[leftmargin=2.6em]
  \item[\textbf{17.}] (1) 最小正周期为 $\pi$；单调递增区间为 $\left[-\dfrac\pi3+k\pi,\dfrac\pi6+k\pi\right]$，$k\in\mathbb Z$；(2) 将 $y=\sin2x$ 的图象向左平移 $\dfrac\pi{12}$，再向上平移 $\dfrac32$
  \item[\textbf{18.}] (1) $AO\perp$ 平面 $BCD$；(2) $\arccos\dfrac{\sqrt2}{4}$；(3) $\dfrac{\sqrt{21}}7$
  \item[\textbf{19.}] (1) $17.5$ 升；(2) 速度为 $80$ 千米/小时，最少耗油 $11.25$ 升
  \item[\textbf{20.}] (1) $\left(x+\dfrac12\right)^2+(y-\sqrt2)^2=\dfrac94$ 或 $\left(x+\dfrac12\right)^2+(y+\sqrt2)^2=\dfrac94$；(2) $-\dfrac12<x_G<0$
  \item[\textbf{21.}] (1) $h(t)=\begin{cases}-t^2+6t+7,&t\le3,\\16,&3<t<4,\\-t^2+8t,&t\ge4.\end{cases}$；(2) 存在，且 $7<m<15-6\ln3$
  \item[\textbf{22.}] (1) $a_n=2^n-1$；(2) $\{b_n\}$ 是等差数列；具体为 $b_n=2+(n-1)(b_2-2)$，其中 $b_2$ 可任取实数；(3) $\dfrac n2-\dfrac13<\displaystyle\sum_{k=1}^n\dfrac{a_k}{a_{k+1}}<\dfrac n2$
\end{itemize}
